% from ega3-1-fr.tex
\oldpage[0\textsubscript{III}]{26}
l'hypothèse que E et $E'$ sont régulières entraîne aussi
$u_2^{pq}(Z_\infty(E_2^{pq}))=Z_\infty(E_2^{\prime pq})$, et comme
$u_2^{pq}$ est un isomorphisme, il en est de même de $u_\infty^{pq}$~; on en
conclut donc que $\operatorname{gr}_p(u^{p+q})$ est aussi un isomorphisme.
Mais comme les filtrations des $E^n$ et des $E^{\prime n}$ sont discrètes et
exhaustives, cela entraîne que les $u^n$ sont aussi des isomorphismes
(Bourbaki, \emph{Alg. comm.}, chap. III, \S~2, n\textsuperscript{o} 8,
th. 1).
\end{proof}

\begin{env}[11.1.6]
\label{0.11.1.6-fr}
Il résulte de (\hyperref[0.11.1.1.2-fr]{11.1.1.2}) et de la définition
(\hyperref[0.11.1.1.3-fr]{11.1.1.3}) que si, dans une suite spectrale E, on a
$E_r^{pq}=0$, on a $E_k^{pq}=0$ pour $k\geq r$ et $E_\infty^{pq}=0$. On dit
qu'une suite spectrale est \emph{dégénérée} s'il existe un entier $r\geq2$ et,
pour tout entier $n\in\mathbb Z$, un entier $q(n)$ tel que
$E_r^{n-q,q}=0$ pour tout $q\ne q(n)$. On déduit d'abord de la remarque
précédente que l'on a aussi $E_k^{n-q,q}=0$ pour $k\geq r$ (y compris
$k=\infty$) et $q\ne q(n)$. En outre, la définition de $E_{r+1}^{pq}$ montre
que l'on a
$E_{r+1}^{n-q(n),q(n)}=E_r^{n-q(n),q(n)}$~; si E est faiblement convergente,
on a donc aussi
$E_\infty^{n-q(n),q(n)}=E_r^{n-q(n),q(n)}$~; autrement dit, pour tout
$n\in\mathbb Z$, $\operatorname{gr}_p(E^n)=0$ pour $p\ne q(n)$ et
$\operatorname{gr}_{q(n)}(E^n)=E_r^{n-q(n),q(n)}$. Si en outre la filtration
de $E^n$ est discrète et exhaustive, la suite E est régulière et on a
$E^n=E_r^{n-q(n),q(n)}$ à un isomorphisme près.
\end{env}

\begin{env}[11.1.7]
\label{0.11.1.7-fr}
Supposons que dans la catégorie $\mathcal C$ les limites inductives filtrantes
existent et soient exactes, et soit $(E_\lambda,u_{\mu\lambda})$ un système
inductif (suivant un ensemble d'indices filtrant) de suites spectrales de
$\mathcal C$. Alors la \emph{limite inductive} de ce système inductif existe
dans la catégorie additive des suites spectrales d'objets de $\mathcal C$~:
il suffit pour le voir de définir $E_r^{pq}$, $d_r^{pq}$, $\alpha_r^{pq}$,
$B_\infty(E_2^{pq})$, $Z_\infty(E_2^{pq})$, $E^n$, $F^p(E^n)$ et
$\beta^{pq}$ comme limites inductives respectives de $E_{r,\lambda}^{pq}$,
$d_{r,\lambda}^{pq}$, $\alpha_{r,\lambda}^{pq}$,
$B_\infty(E_{2,\lambda}^{pq})$, $Z_\infty(E_{2,\lambda}^{pq})$,
$E_\lambda^n$, $F^p(E_\lambda^n)$ et $\beta_\lambda^{pq}$~; la vérification
des conditions de (\hyperref[0.11.1.1-fr]{11.1.1}) résulte de l'exactitude du
foncteur $\varinjlim$ dans $\mathcal C$.
\end{env}

\begin{remark}[11.1.8]
\label{0.11.1.8-fr}
Supposons que la catégorie $\mathcal C$ soit la catégorie des $A$-modules sur
un anneau \emph{noethérien} $A$ (resp. sur un anneau $A$). Alors, les
définitions de (\hyperref[0.11.1.1-fr]{11.1.1}) montrent que si pour un $r$
donné, les $E_r^{pq}$ sont des $A$-modules \emph{de type fini} (resp. de
\emph{longueur finie}), il en est de même de tous les modules $E_s^{pq}$ pour
$s\geq r$, ainsi que des $E_\infty^{pq}$. Si en outre la filtration de
l'aboutissement $(E^n)$ est discrète et co-discrète pour tout $n$, on en
conclut que chacun des $E^n$ est aussi un $A$-module \emph{de type fini}
(resp. de \emph{longueur finie}).
\end{remark}

\begin{env}[11.1.9]
\label{0.11.1.9-fr}
Nous aurons à considérer des conditions assurant qu'une suite spectrale E est
birégulière de façon «~uniforme~» en $p+q=n$. On utilisera alors le lemme
suivant~:
\end{env}

\begin{lemma}[11.1.10]
\label{0.11.1.10-fr}
Soit $(E_r^{pq})$ une famille d'objets de $\mathcal C$ liés par les données
a), b), c) de (\hyperref[0.11.1.1-fr]{11.1.1}). Pour un entier fixé $n$, les
propriétés suivantes sont équivalentes~:
\begin{enumerate}
  \item[(a)] Il existe un entier $r(n)$ tel que pour $r\geq r(n)$,
    $p+q=n$ ou $p+q=n-1$, les morphismes $d_r^{pq}$ soient tous nuls.
  \item[(b)] Il existe un entier $r(n)$ tel que pour $p+q=n$ ou $p+q=n+1$,
    on ait $B_r(E_2^{pq})=B_s(E_2^{pq})$ pour $s\geq r\geq r(n)$.
  \item[(c)] Il existe un entier $r(n)$ tel que pour $p+q=n$ ou $p+q=n-1$,
    on ait $Z_r(E_2^{pq})=Z_s(E_2^{pq})$ pour $s\geq r\geq r(n)$.
  \item[(d)] Il existe un entier $r(n)$ tel que pour $p+q=n$, on ait
    $B_r(E_2^{pq})=B_s(E_2^{pq})$ et $Z_r(E_2^{pq})=Z_s(E_2^{pq})$ pour
    $s\geq r\geq r(n)$.
\end{enumerate}
\end{lemma}

\begin{proof}
